\section{Заключение}
    
% ============================================================
% ЗАКЛЮЧЕНИЕ (2--3 страницы)
% ============================================================

% ЦЕЛЬ: Подвести итоги всей работы.
% ЧТО ОПИСАТЬ:
%   - Перечисление решённых задач (по пунктам из Введения)
%   - Основные результаты: работающая система, 
%     подтверждённое качество распознавания
%   - Практическая ценность: готовый продукт 
%     для локальной транскрибации
%   - Направления развития:
%     — GPU-ускорение
%     — Магазин моделей
%     — Предобработка аудио (шумоподавление)
%     — Интеграция LLM для анализа текста
%     — Расширение GUI до редактора


% ============================================================
% СПИСОК ЛИТЕРАТУРЫ (≥25 источников)
% ============================================================
\section*{Список литературы}

% Категории источников:
%   - Статьи по ASR: Whisper (Radford et al., 2023), 
%     Conformer, CTC/Attention
%   - Статьи по диаризации: pyannote (Bredin et al.), EEND
%   - Статьи по LID: ECAPA-TDNN, VoxLingua107
%   - Паттерны проектирования: GoF, Clean Architecture (Martin)
%   - ГОСТы: ГОСТ 19.201, ГОСТ 34.602
%   - Документация: pyannote, speechbrain, HuggingFace, 
%     faster-whisper, PyQt6